\begin{convention}[Index suppression and aggregation from the root variable]
\label{conv:notation}

Every variable in the model is defined with a given root indices. For example, variable $\phi$, is defined with the following indices $(c\in C,i\in V,j\in V,m\in \mathcal{M})$; we call this the \textbf{root variable}.
Writing a variable with fewer indices always denotes a
\emph{derived variables}: one or more indices have been suppressed and
the root values aggregated according to the rules below.



% ------------------------------------------------------------------
\medskip
\noindent\textbf{Rule~1 (Sum aggregation).
 }

Dropping any indeces from any variable $\phi$ stands for the sum of
the root variables over every value of those indices:
%k
\begin{align*}
  \phi_{cij}
    &\;\coloneqq\; \sum_{m\in\mathcal{M}} \phi_{cijm},
  &
  \phi_{ijm}
    &\;\coloneqq\; \sum_{c\in C}           \phi_{cijm},
  &
  \phi
    &\;\coloneqq\;
      \sum_{c\in C}\sum_{i\in V}\sum_{j\in V}\sum_{m\in\mathcal{M}} \phi_{cijm}.
\end{align*}
%
For continuous and integer variables this is the natural
interpretation.
For binary variables the result is an \emph{integer count} of how
many root-level binaries in the group equal~1; it is not itself
binary.
For example, $X_{cij} \coloneqq \sum_{m\in\mathcal{M}}
X_{cijm}$ counts how many modes are active for type $c$ on
lane $(i,j)$ and takes values in $\{0,1,\dots,|\mathcal{M}|\}$.
The sum aggreagtion can be considered a shorthand notation, since we do not need to create new variables in the solver. We can also create derived $\phi$ variable to the model and the associated constraint.:while

% ------------------------------------------------------------------
\medskip
\noindent\textbf{Rule~2 (Dot notation for the node-pair $(i,j)$).}

The root variable carries two active indices from the same set $V$.
Suppressing one while keeping the other would be ambiguous, so a
dot~$(\,\cdot\,)$ marks the suppressed position explicitly:
%
\begin{align*}
  \phi_{c\cdot j m}
    &\;\coloneqq\; \sum_{i\in V} \phi_{cijm},
  &
  \phi_{ci \cdot m}
    &\;\coloneqq\; \sum_{j\in V} \phi_{cijm}.
\end{align*}
%
The dot is only ever needed for the origin/destination pair, as no
other pair of active indices belongs to the same set.
Rule~1 applies as normal: the result is a sum, and for binary
variables it is an integer count.

% ------------------------------------------------------------------
\medskip
\noindent\textbf{Rule~3 (OR-aggregated binaries).
  }

We can also aggregate binary variables to denote an OR aggreagtion. Then, a new binary decision variable is introduced,
denoted with an overline ($\overline{X}$).
This variable is not a notational shorthand: it must be declared
in the solver and its relationship to the root-level binaries
enforced by an explicit big-$M$ pair, or boolean imposition.
Crucially, the pair is written directly in terms of the
sum aggregate defined by Rule~1, so the sum-shorthand appears
in the constraints themselves.

\smallskip
\noindent\emph{Suppressing the mode index $m$} yields the new
binary $\overline{X}_{cij}$, linked to the sum aggregate
$X_{cij} = \sum_{m\in\mathcal{M}} X_{cijm}$ by:
%
\begin{align}
  \overline{X}_{cij}
    &\leq X_{cij}
  && \forall\,(i,j)\in E,\ c\in C
  \label{eq:or-m-lb}\\
  X_{cij}
    &\leq |M|\cdot\overline{X}_{cij}
  && \forall\,(i,j)\in E,\ c\in C
  \label{eq:or-m-ub}
\end{align}
%
Constraint~\eqref{eq:or-m-lb} forces $\overline{X}_{cij}$
to~0 when the count $X_{cij}$ is~0 (no mode active).
Constraint~\eqref{eq:or-m-ub} forces $\overline{X}_{cij}$
to~1 when the count is positive (at least one mode active).
Together they enforce
$\overline{X}_{cij} = 1 \iff X_{cij} \geq 1$.

% ------------------------------------------------------------------
\medskip
\noindent\textbf{Rule~4 (At Least One)}

We can also aggregate binary variables to denote an At Least One (ALO) aggreagtion. Then, a new binary decision variable is introduced,
denoted with an overline  $\overline{X}_{cij}$. We also need to
%
\begin{align}
  X_{cij}
    &= 1
  && \forall\,(i,j)\in E
  \label{eq:exactly-one}
\end{align}
%


% ------------------------------------------------------------------
%  TIKZ LATTICE
% ------------------------------------------------------------------
\begin{figure}[H]
\centering
\resizebox{0.95\linewidth}{!}{%
  % --- ADJUSTABLE PARAMETERS ---
\def\hsep{4.0}
\def\vsep{1.4}
\def\logicXshift{7.5}
% -----------------------------
\begin{tikzpicture}[
  every node/.style = {
    draw, rounded corners=3pt, inner sep=5pt,
    align=center, font=\small, thick
  },
  rootBox/.style   = {fill=blue!10, draw=blue!60!black},
  sumBox/.style    = {fill=orange!10, draw=orange!60!black},
  binBox/.style    = {fill=red!5, draw=red!80!black, dashed},
  arr/.style       = {->, >=stealth, thick},
  barr/.style      = {->, >=stealth, thick, dashed, red!80!black},
  lbl/.style       = {draw=none, fill=none, font=\footnotesize\itshape, inner sep=2pt},
  blbl/.style      = {draw=none, fill=none, font=\footnotesize\bfseries, inner sep=2pt, text=red!80!black},
  graylbl/.style   = {draw=none, fill=none, font=\small\bfseries, gray!50}
]

%% --- SUM AGGREGATION SECTION (Left) ---
\node[rootBox] (root) at (0,0) {$\phi_{cijm}$};
\node[sumBox] (Fcij)  at (-\hsep/2, \vsep)
  {$\widetilde{\phi}_{cij} \coloneqq \sum_{m \in \mathcal{M}} \phi_{cijm}$};
\node[sumBox] (Fijm)  at (\hsep/2, \vsep)
  {$\widetilde{\phi}_{ijm} \coloneqq \sum_{c \in C} \phi_{cijm}$};
\node[sumBox] (Fij)   at (0, 2*\vsep)
  {$\widetilde{\phi}_{ij} \coloneqq \sum_{c \in C} \sum_{m \in \mathcal{M}} \phi_{cijm}$};

\draw[arr] (root.north) -- (Fcij.south);
\draw[arr] (root.north) -- (Fijm.south);
\draw[arr] (Fcij.north) -- (Fij.south);
\draw[arr] (Fijm.north) -- (Fij.south);

% Dot notation
\node[sumBox] (dotj) at (-\hsep/2, -\vsep)
  {$\widetilde{\phi}_{c \cdot j m} \coloneqq \sum_{i \in V} \phi_{cijm}$};
\node[sumBox] (doti) at (\hsep/2, -\vsep)
  {$\widetilde{\phi}_{ci \cdot m} \coloneqq \sum_{j \in V} \phi_{cijm}$};

\draw[arr] (root.south) -- (dotj.north);
\draw[arr] (root.south) -- (doti.north);

%% --- BINARY LOGIC AGGREGATION SECTION (Right) ---
\begin{scope}[xshift=\logicXshift cm]
\node[rootBox] (Xroot) at (0, 0) {$X_{cijm}$};
\node[binBox] (XOR)  at (0, \vsep*1.36)
  {$\widecheck{X}_{cij}$ (OR aggregate) \\
   \scriptsize $\widecheck{X}_{cij} \leq \widetilde{X}_{cij}$ \\
   \scriptsize $\widetilde{X}_{cij} \leq |\mathcal{M}|\,\widecheck{X}_{cij}$};
\node[binBox] (XHAT) at (0, -\vsep*1.36)
  {$\widehat{X}_{cij}$ (At-most-one) \\
   \scriptsize $\widehat{X}_{cij} \leq \widetilde{X}_{cij}$ \\
   \scriptsize $\widetilde{X}_{cij} \leq \widehat{X}_{cij}$};

\draw[barr] (Xroot.north) -- (XOR.south)  node[blbl, midway, right] {Rule 3};
\draw[barr] (Xroot.south) -- (XHAT.north) node[blbl, midway, right] {Rule 4};
\end{scope}

%% --- DIVIDER AND LABELS ---
\draw[thick, gray!30, dashed]
  (\logicXshift/2, -1.5*\vsep) -- (\logicXshift/2, 2.8*\vsep);
\node[graylbl] at (0,            2.7*\vsep) {Sum Aggregation};
\node[graylbl] at (\logicXshift, 2.7*\vsep) {Binary Logic Aggregation};

\end{tikzpicture}
}
\caption{%
  Derivation lattice for a generic variable $\phi$ (left panel,
  continuous or integer) and its binary counterpart $X$ (right
  panel), separated by the dashed vertical rule.
  On the left, every arrow represents a sum aggregation
  (Rule~1--2): purely notational, no new variable introduced.
  On the right, solid arrows represent the same sum aggregation
  applied to binary variables, producing an integer count.
  Dashed arrows introduce a genuinely new solver variable (dashed
  border): the OR aggregate $\overline{X}$ (Rule~3), defined as a
  binary proxy for whether the count is positive, enforced by the
  big-$M$ pair shown in each node.
  The hat variant $\hat{X}$ (Rule~4) is the same variable as
  $\overline{X}_{cij}$ with the additional
  constraint~\eqref{eq:exactly-one} in force.%
}
\label{fig:index-lattice}
\end{figure}

\end{convention}
